
\documentclass[a4paper]{article} %
\usepackage{graphicx,amssymb} %

\textwidth=15cm \hoffset=-1.2cm %
\textheight=25cm \voffset=-2cm %

\pagestyle{empty} %

\date{} %

\def\keywords#1{\begin{center}{\bf Keywords}\\{#1}\end{center}} %

\def\titulo#1{\title{#1}} %
\def\autores#1{\author{#1}} %

% Please, do not change any of the above lines








\begin{document}

% Type down your paper title
\titulo{Padovan-like sequences and Bell polynomials}

% Authors
\autores{Nikita Gogin \\ %
      \AA bo Akademi University (Finland) \\ \\ % Affiliation 1
       Aleksandr Myll\"ari \\ % If any other author with different Affilation
       St. George's University (Grenada) \\ \\ % Affiliation 2 (if needed)
       % New author \\
       % New affiliation \\
       % Add authors and affiliation as needed
       \tt{amyllari@sgu.edu} % Only one corresponding e-mail
       }%


\maketitle

\thispagestyle{empty}



% The abstract

\begin{abstract}
We study a class of Padovan-like sequences that can be generated using special matrices of the third order. We show that terms of any sequence of this class can be expressed via Bell polynomials and their derivatives that use as arguments terms of another such sequence with smaller indexes. CAS {\it Mathematica} is used for cumbersome calculations and hypothesis testing.


\end{abstract}

\keywords{Padovan sequence, Fibonacci numbers, Bell polynomials, integer sequences} % Write down at least 3 Keywords






 \section{Introduction}

Integer sequences appear in many branches of science. One famous example is Fibonacci numbers that have been known for more than two thousand years and find applications in mathematics, biology, economics, computer science. Padovan numbers are much younger -� they were introduced only recently \cite{bibl1}.   Below, we will study Padovan-like sequences that can be generated using special matrices of the third order. We will find expressions for terms of one sequence in terms of another sequence via Bell polynomials. CAS {\it Mathematica} was used for cumbersome calculations and hypothesis testing.

Let 
$$A_\alpha = \left( \begin{array}{ccc}
0 & \alpha & 1 \\
1 & 0 & 0 \\
0 & 1 & 0
\end{array} \right) $$ 
and let $\left( \begin{array}{c}
u_n \\ v_n \\ w_n
\end{array} \right)$ for a time  denote the first column of $A_\alpha^n,$  $n\geq 0.$ Then 
\begin{equation} \label{myll7}
u_{n+1} = \alpha v_n + w_n, \quad v_{n+1} = u_n, \quad w_{n+1} = v_n.
\end{equation}
%Let $u_0=\alpha, v_0=0, w_0=1.$ 
We have
\begin{equation} \label{myll8}
u_{n+1} = \alpha u_{n-1} + u_{n-2}, \quad u_0 =1, \quad u_1= 0, \quad u_2=\alpha, \quad (u_{-1}= 0,)
\end{equation}
  
\begin{equation} \label{myll9}
A_\alpha^n = \left( \begin{array}{ccc}
u_n & u_{n+1} & u_{n-1} \\
u_{n-1} & u_n & u_{n-2} \\
u_{n-2} & u_{n-1} & u_{n-3}
\end{array} \right). 
\end{equation}

Some examples of sequences generated by the matrix $A_\alpha$ with references to the On-Line Encyclopedia of Integer Sequences (OEIS, http://oeis.org/) are given in Table 1.

\begin{table}[h!]
\caption{Examples of sequences generated by the matrix $A_\alpha.$}
\begin{center}
\begin{tabular}{|l|l|l|l|}
\hline &&& \\
$\alpha$ & OEIS  & First terms & Comment \\ & reference && \\ \hline  &&& \\
 $\alpha = 1$ &
A000931 &	1,0,1,1,1,2,2,3,4,5,7,9,12,16,21, & Padovan sequence : $u_n = p_n$  \\&&& \\ \hline &&& \\
  $\alpha = 2$ &
A008346 &	1,0,2,1,4,4,9,12,22,33,56,88,145, & $u_n = f_n = Fibonacci(n) + (-1)^n $ \\ &&& \\ \hline &&& \\
  $\alpha = 3$ &
A052931  &	1,0,3,1,9,6,28,27,90,109,297,417, &
\\ &&& \\ \hline &&&\\
   $\alpha = 0$ &
 A079978&	1,0,0,1,0,0,1,0,0,1,0,0,1,0,0, 1,  & $ \begin{array}{l} u_n = t_n = \left\{
\begin{array}{ll}
1, & n \equiv 0(mod \: 3) \\ 
0, & n \equiv 1, 2 (mod \: 3)
\end{array} \right. \\
 = {\frac n 3 \choose 0} \end{array} $\\ &&& \\ \hline &&& \\
   $\alpha = -1$ &
A077961&	1,0,-1,1,1,-2,0,3,-2,-3,5,1,-8,4,9, &
\\ &&& \\ \hline &&& \\
   $\alpha = -2$ &
A077965 &	1,0,-2,1,4,-4,-7,12,10,-31,-8,72,-15,-152,102 &
\\ &&& \\
\hline
\end{tabular}
\end{center}
\end{table}




 \section{Main Result}
	Now let $\beta \in Z,$ and let $v_n$ denote terms of the sequence corresponding to the powers of the matrix $A_\beta,$ and let $w_n$ denote terms of the sequence corresponding to the powers of the matrix $A_{\alpha + \beta}:$
$$
A_{\alpha + \beta}^n = \left( \begin{array}{ccc}
w_n & w_{n+1} & w_{n-1} \\
w_{n-1} & w_n & w_{n-2} \\
w_{n-2} & w_{n-1} & w_{n-3}
\end{array} \right)
= \left( \begin{array}{ccc}
0 & \alpha + \beta & 1 \\
1 & 0 & 0 \\
0 & 1 & 0
\end{array} \right)^n = (A_\alpha + \beta e)^n = (A_\beta + \alpha e)^n,
$$
where  $e= \left(
\begin{array}{ccc}
0 & 1 & 0 \\
0 & 0 & 0 \\
0 & 0 & 0
\end{array}
\right). $


Terms of the sequence generated by  $A_{\alpha + \beta}$ can be expressed in terms of the sequence generated by  $A_\alpha$ as follows:
\begin{eqnarray} \label{myll12}
w_n &=& \sum_{\varepsilon = 0}^{1} \sum_{s = 1}^{\left[ \frac{n + \varepsilon}2 \right]} \frac{(s-\varepsilon)! \cdot \beta^{s -\varepsilon}}{(n-s+\varepsilon)!} \times \\ \nonumber
&& \times \mathcal{D}^{\varepsilon}( B_{n-s + \varepsilon, s} 
 (1! \cdot u_{k_1 -1}, 2! \cdot  u_{k_2 -1}, \dots , (n-2s +1 +\varepsilon)! \cdot u_{k_{n-2s +1 +\varepsilon}-1})),
\end{eqnarray}
where ${k_i}$ run over all partitions  of $n-s+\varepsilon$ into $s$ parts, $\mathcal{D}^0 = id$ - identity operator, $\mathcal{D}^1 = \mathcal{D} = \sum_{r=1}^{(\infty)} x_{r+1} \frac{\partial}{\partial x_r},$ and
 \begin{equation} \label{myll5}  
 B_{n,k} (x_1, x_2, \dots, x_{n-k+1}) = \sum \frac{n!}{j_1!j_2! \dots j_{n-k+1}!} \left( \frac{x_1}{1!} \right)^{j_1}  \left( \frac{x_2}{2!} \right)^{j_2} \dots  \left( \frac{x_{n-k+1}}{(n-k+1)!} \right)^{j_{n-k+1}} 
\end{equation} 
are partial Bell polynomials, and summing is done for all sets $j_1, j_2, \dots, j_{n-k+1}$ of non-negative integers such that $j_1 + j_2 +  \dots + j_{n-k+1} = k$ and $1 \cdot j_1 + 2 \cdot j_2 + 3 \cdot j_3 +  \dots + (n- k +1) \cdot j_{n-k+1} = n$
 (see \cite{bibl3}.)
     

\section{Examples}
\subsection{Example 1}
          For $\alpha = \beta =1$   formula (\ref{myll12}) gives a relation between terms of Fibonacci and Padovan sequences:
\begin{eqnarray}\label{myll13}
f_n &=& Fibonacci(n) +(-1)^n =
\\
&=& \sum_{\varepsilon = 0}^1 \sum_{s=1}^{\left[ \frac{n + \varepsilon} 2  \right]} \frac{(s - \varepsilon)!}{(n - s + \varepsilon)!} \mathcal{D}^{\varepsilon} ( B_{n-s+\varepsilon,s}  (1! \cdot p_{k_1 -1}, 2! \cdot  p_{k_2 -1}, \dots , (n-2s +1 +\varepsilon)! \cdot p_{k_{n-2s +1 +\varepsilon}-1})),
\nonumber 
\end{eqnarray}
and thus for $n \geq 1$ the $n-$th term of the sequence $Fibonacci(n) + (-1)^n$   is expressed as sums of products of the $n$ first terms of the Padovan sequence. 

\subsection{Example 2}
 For $\alpha =2$  and $\beta = -2$  formula (4) gives a relation between terms of sequences $f_n =Fibonacci(n) + (-1)^n$ and $t_n =\left\{ 
\begin{array}{ll}
1, & n \equiv 0 \; (mod \; 3) \\
0, & n \equiv 1, 2 \; (mod \; 3)
\end{array}  
\right. ,$
$n \geq 1:$
\begin{eqnarray}\nonumber
\sum_{\varepsilon = 0}^1 \sum_{s=1}^{\left[ \frac{n + \varepsilon} 2  \right]} \frac{(s - \varepsilon)! (-2)^{s- \varepsilon}}{(n - s + \varepsilon)!} \mathcal{D}^{\varepsilon} ( B_{n-s+\varepsilon,s}  (1! \cdot f_{k_1 -1}, 2! \cdot  f_{k_2 -1}, \dots , (n-2s +1 +\varepsilon)! \cdot f_{k_{n-2s +1 +\varepsilon}-1}))= \\
= \left\{
\begin{array}{ll}
1, & n \equiv 0 \; (mod \; 3) \\
0, & n \equiv 1, 2 \; (mod \; 3)
\end{array}  
\right.
\label{myll14}
\end{eqnarray}

\subsection{Example 3}

 In particular, the terms of every Padovan-like sequence with $\alpha \neq 0 $  can be expressed in the same way via the terms of  the "simplest generator"  of the class, namely the sequence $1,0,0,1,0,0,1,0,0,...$  generated by $\alpha = 0$ (OEIS-number  A079978, see Table 1). This statement is a far-reaching generalization of the result announced by J. Vladetta in \cite{bibl6}.


\begin{thebibliography}{9}
\bibitem{bibl1}  Richard Padovan, Dom Hans Van Der Laan and the Plastic Number, pp. 181-193 in Nexus IV: \textit{Architecture and Mathematics}, eds. Kim Williams and Jose Francisco Rodrigues Fucecchio (Florence): Kim Williams Books, 2002.



 %George E. Andrews, \textit{The Theory of Partitions }  Add.-Wesley PC, 1976.

%\bibitem{bibl2} http://en.wikipedia.org/wiki/Composition$\_$(number$\_$theory) 

\bibitem{bibl3} http://en.wikipedia.org/wiki/Bell$\_$polynomials

%\bibitem{bibl4} Paul Barry \textit{} in OEIS

%\bibitem{bibl5} http://en.wikipedia.org/wiki/Padovan$\_$polynomials

\bibitem{bibl6} Vladetta.J \textit{}, in "Sloane's A000931 : Padovan sequence", The On-Line Encyclopedia of Integer Sequences. OEIS Foundation. 
 

\end{thebibliography}

\end{document}
